\section{Completion for Noetherian rings}
\label{section-completion-noetherian}

\noindent
In this section we discuss completion with respect to ideals in
Noetherian rings.

\begin{lemma}
\label{lemma-completion-tensor}
Let $I$ be an ideal of a Noetherian ring $R$.
Denote ${}^\wedge$ completion with respect to $I$.
\begin{enumerate}
\item If $K \to N$ is an injective map of finite $R$-modules,
then the map on completions $K^\wedge \to N^\wedge$ is injective.
\item If $0 \to K \to N \to M \to 0$ is a short exact sequence
of finite $R$-modules, then $0 \to K^\wedge \to N^\wedge \to M^\wedge \to 0$
is a short exact sequence.
\item If $M$ is a finite $R$-module, then $M^\wedge = M \otimes_R R^\wedge$.
\end{enumerate}
\end{lemma}

\begin{proof}
Setting $M = N/K$ we find that part (1) follows from part (2).
Let $0 \to K \to N \to M \to 0$ be as in (2).
For each $n$ we get the short exact sequence
$$
0 \to K/(I^nN \cap K) \to N/I^nN \to M/I^nM \to 0.
$$
By Lemma \ref{lemma-Mittag-Leffler}
we obtain the exact sequence
$$
0 \to \lim K/(I^nN \cap K) \to N^\wedge \to M^\wedge \to 0.
$$
By the Artin-Rees Lemma \ref{lemma-Artin-Rees} we may choose $c$ such that
$I^nK \subset I^n N \cap K \subset I^{n-c} K$ for $n \geq c$.
Hence $K^\wedge = \lim K/I^nK = \lim K/(I^nN \cap K)$
and we conclude that (2) is true.

\medskip\noindent
Let $M$ be as in (3) and let $0 \to K \to R^{\oplus t} \to M \to 0$
be a presentation of $M$. We get a commutative diagram
$$
\xymatrix{
&
K \otimes_R R^\wedge \ar[r] \ar[d] &
R^{\oplus t} \otimes_R R^\wedge \ar[r] \ar[d] &
M \otimes_R R^\wedge \ar[r] \ar[d] &
0 \\
0 \ar[r] &
K^\wedge \ar[r] &
(R^{\oplus t})^\wedge \ar[r] &
M^\wedge \ar[r] & 0
}
$$
The top row is exact, see Section \ref{section-flat}.
The bottom row is exact by part (2).
By Lemma \ref{lemma-completion-generalities}
the vertical arrows are surjective.
The middle vertical arrow is an isomorphism.
We conclude (3) holds by the Snake Lemma \ref{lemma-snake}.
\end{proof}

\begin{lemma}
\label{lemma-completion-flat}
Let $I$ be an ideal of a Noetherian ring $R$.
Denote ${}^\wedge$ completion with respect to $I$.
\begin{enumerate}
\item The ring map $R \to R^\wedge$ is flat.
\item The functor $M \mapsto M^\wedge$ is exact on the category of
finitely generated $R$-modules.
\end{enumerate}
\end{lemma}

\begin{proof}
Consider $J \otimes_R R^\wedge \to R \otimes_R R^\wedge = R^\wedge$
where $J$ is an arbitrary ideal of $R$.
According to Lemma \ref{lemma-completion-tensor} this
is identified with $J^\wedge \to R^\wedge$ and $J^\wedge \to R^\wedge$
is injective. Part (1) follows from Lemma \ref{lemma-flat}.
Part (2) is a reformulation of
Lemma \ref{lemma-completion-tensor} part (2).
\end{proof}

\begin{lemma}
\label{lemma-completion-faithfully-flat}
Let $I$ be an ideal of a Noetherian ring $R$. Denote $R^\wedge$
the completion of $R$ with respect to $I$. If $I$ is contained
in the Jacobson radical of $R$, then the ring map $R \to R^\wedge$
is faithfully flat. In particular, if $(R, \mathfrak m)$ is a Noetherian
local ring, then the completion $\lim_n R/\mathfrak m^n$ is faithfully flat.
\end{lemma}

\begin{proof}
By Lemma \ref{lemma-completion-flat} it is flat.
The composition $R \to R^\wedge \to R/I$ where
the last map is the projection map $R^\wedge \to R/I$ shows that
any maximal ideal of $R$ is in the image of $\Spec(R^\wedge)
\to \Spec(R)$. Hence the map is faithfully
flat by Lemma \ref{lemma-ff}.
\end{proof}

\begin{lemma}
\label{lemma-completion-complete}
Let $R$ be a Noetherian ring.
Let $I$ be an ideal of $R$.
Let $M$ be an $R$-module.
Then the completion $M^\wedge$
of $M$ with respect to $I$ is $I$-adically complete,
$I^n M^\wedge = (I^nM)^\wedge$, and $M^\wedge/I^nM^\wedge = M/I^nM$.
\end{lemma}

\begin{proof}
This is a special case of
Lemma \ref{lemma-hathat-finitely-generated}
because $I$ is a finitely generated ideal.
\end{proof}

\begin{lemma}
\label{lemma-completion-Noetherian}
Let $I$ be an ideal of a ring $R$. Assume
\begin{enumerate}
\item $R/I$ is a Noetherian ring,
\item $I$ is finitely generated.
\end{enumerate}
Then the completion $R^\wedge$ of $R$ with respect to $I$
is a Noetherian ring complete with respect to $IR^\wedge$.
\end{lemma}

\begin{proof}
By
Lemma \ref{lemma-hathat-finitely-generated}
we see that $R^\wedge$ is $I$-adically complete. Hence it is also
$IR^\wedge$-adically complete. Since $R^\wedge/IR^\wedge = R/I$ is
Noetherian we see that after replacing $R$ by $R^\wedge$ we may in
addition to assumptions (1) and (2) assume that also $R$ is $I$-adically
complete.

\medskip\noindent
Let $f_1, \ldots, f_t$ be generators of $I$.
Then there is a surjection of rings
$R/I[T_1, \ldots, T_t] \to \bigoplus I^n/I^{n + 1}$
mapping $T_i$ to the element $\overline{f}_i \in I/I^2$.
Hence $\bigoplus I^n/I^{n + 1}$ is a Noetherian ring.
Let $J \subset R$ be an ideal. Consider the ideal
$$
\bigoplus J \cap I^n/J \cap I^{n + 1} \subset \bigoplus I^n/I^{n + 1}.
$$
Let $\overline{g}_1, \ldots, \overline{g}_m$ be generators of this
ideal. We may choose $\overline{g}_j$ to be a homogeneous element
of degree $d_j$ and we may pick $g_j \in J \cap I^{d_j}$ mapping to
$\overline{g}_j \in J \cap I^{d_j}/J \cap I^{d_j + 1}$. We claim
that $g_1, \ldots, g_m$ generate $J$.

\medskip\noindent
Let $x \in J \cap I^n$. There exist $a_j \in I^{\max(0, n - d_j)}$ such that
$x - \sum a_j g_j \in J \cap I^{n + 1}$.
The reason is that $J \cap I^n/J \cap I^{n + 1}$ is equal to
$\sum \overline{g}_j I^{n - d_j}/I^{n - d_j + 1}$ by our choice
of $g_1, \ldots, g_m$. Hence starting with $x \in J$ we can find
a sequence of vectors $(a_{1, n}, \ldots, a_{m, n})_{n \geq 0}$
with $a_{j, n} \in I^{\max(0, n - d_j)}$ such that
$$
x =
\sum\nolimits_{n = 0, \ldots, N}
\sum\nolimits_{j = 1, \ldots, m} a_{j, n} g_j \bmod I^{N + 1}
$$
Setting $A_j = \sum_{n \geq 0} a_{j, n}$ we see that
$x = \sum A_j g_j$ as $R$ is complete. Hence $J$ is finitely generated and
we win.
\end{proof}

\begin{lemma}
\label{lemma-completion-Noetherian-Noetherian}
Let $R$ be a Noetherian ring.
Let $I$ be an ideal of $R$.
The completion $R^\wedge$ of $R$ with respect to $I$ is
Noetherian.
\end{lemma}

\begin{proof}
This is a consequence of
Lemma \ref{lemma-completion-Noetherian}.
It can also be seen directly as follows.
Choose generators $f_1, \ldots, f_n$ of $I$.
Consider the map
$$
R[[x_1, \ldots, x_n]] \longrightarrow R^\wedge,
\quad
x_i \longmapsto f_i.
$$
This is a well defined and surjective ring map
(details omitted).
Since $R[[x_1, \ldots, x_n]]$ is Noetherian (see
Lemma \ref{lemma-Noetherian-power-series}) we win.
\end{proof}

\noindent
Suppose $R \to S$ is a local homomorphism of local rings $(R, \mathfrak m)$
and $(S, \mathfrak n)$. Let $S^\wedge$ be the completion
of $S$ with respect to $\mathfrak n$. In general $S^\wedge$ is not
the $\mathfrak m$-adic completion of $S$. If
$\mathfrak n^t \subset \mathfrak mS$ for some $t \geq 1$
then we do have $S^\wedge = \lim S/\mathfrak m^nS$ by
Lemma \ref{lemma-change-ideal-completion}. In some cases this even
implies that $S^\wedge$ is finite over $R^\wedge$.

\begin{lemma}
\label{lemma-finite-after-completion}
Let $R \to S$ be a local homomorphism of local rings $(R, \mathfrak m)$
and $(S, \mathfrak n)$. Let $R^\wedge$, resp.\ $S^\wedge$ be the completion
of $R$, resp.\ $S$ with respect to $\mathfrak m$, resp.\ $\mathfrak n$.
If $\mathfrak m$ and $\mathfrak n$ are finitely generated and
$\dim_{\kappa(\mathfrak m)} S/\mathfrak mS < \infty$, then
\begin{enumerate}
\item $S^\wedge$ is equal to the $\mathfrak m$-adic completion of $S$, and
\item $S^\wedge$ is a finite $R^\wedge$-module.
\end{enumerate}
\end{lemma}

\begin{proof}
We have $\mathfrak mS \subset \mathfrak n$ because $R \to S$ is a local
ring map.
The assumption $\dim_{\kappa(\mathfrak m)} S/\mathfrak mS < \infty$
implies that $S/\mathfrak mS$ is an Artinian ring, see
Lemma \ref{lemma-finite-dimensional-algebra}.
Hence has dimension $0$, see
Lemma \ref{lemma-Noetherian-dimension-0},
hence $\mathfrak n = \sqrt{\mathfrak mS}$.
This and the fact that $\mathfrak n$ is finitely generated
implies that $\mathfrak n^t \subset \mathfrak mS$ for
some $t \geq 1$. By
Lemma \ref{lemma-change-ideal-completion}
we see that $S^\wedge$ can be identified with the $\mathfrak m$-adic
completion of $S$. As $\mathfrak m$ is finitely generated we see from
Lemma \ref{lemma-hathat-finitely-generated}
that $S^\wedge$ and $R^\wedge$ are $\mathfrak m$-adically complete.
At this point we may apply
Lemma \ref{lemma-finite-over-complete-ring}
to $S^\wedge$ as an $R^\wedge$-module to conclude.
\end{proof}

\begin{lemma}
\label{lemma-completion-finite-extension}
Let $R$ be a Noetherian ring. Let $R \to S$ be a finite ring map.
Let $\mathfrak p \subset R$ be a prime and let
$\mathfrak q_1, \ldots, \mathfrak q_m$ be the primes of $S$
lying over $\mathfrak p$
(Lemma \ref{lemma-finite-finite-fibres}).
Then
$$
R_\mathfrak p^\wedge \otimes_R S =
(S_\mathfrak p)^\wedge =
S_{\mathfrak q_1}^\wedge \times \ldots \times S_{\mathfrak q_m}^\wedge
$$
where the $(S_\mathfrak p)^\wedge$ is the completion with respect to
$\mathfrak p$ and the local rings $R_\mathfrak p$ and
$S_{\mathfrak q_i}$ are completed with respect to their maximal ideals.
\end{lemma}

\begin{proof}
We may replace $R$ by the localization $R_\mathfrak p$ and
$S$ by $S_\mathfrak p = S \otimes_R R_\mathfrak p$.
Hence we may assume that $R$ is a local Noetherian ring and
that $\mathfrak p = \mathfrak m$ is its maximal ideal.
The $\mathfrak q_iS_{\mathfrak q_i}$-adic completion
$S_{\mathfrak q_i}^\wedge$ is equal to the $\mathfrak m$-adic
completion by Lemma \ref{lemma-finite-after-completion}.
For every $n \geq 1$ prime ideals of $S/\mathfrak m^nS$ are in 1-to-1
correspondence with the maximal ideals
$\mathfrak q_1, \ldots, \mathfrak q_m$ of $S$
(by going up for $S$ over $R$, see
Lemma \ref{lemma-integral-going-up}).
Hence
$S/\mathfrak m^nS = \prod S_{\mathfrak q_i}/\mathfrak m^nS_{\mathfrak q_i}$
by Lemma \ref{lemma-artinian-finite-length}
(using for example
Proposition \ref{proposition-dimension-zero-ring}
to see that $S/\mathfrak m^nS$ is Artinian).
Hence the $\mathfrak m$-adic completion $S^\wedge$ of $S$ is equal to
$\prod S_{\mathfrak q_i}^\wedge$. Finally, we have
$R^\wedge \otimes_R S = S^\wedge$ by Lemma \ref{lemma-completion-tensor}.
\end{proof}

\begin{lemma}
\label{lemma-split-completed-sequence}
Let $R$ be a ring. Let $I \subset R$ be an ideal.
Let $0 \to K \to P \to M \to 0$ be a short exact sequence of
$R$-modules. If $M$ is flat over $R$ and $M/IM$ is a projective
$R/I$-module, then the sequence of $I$-adic completions
$$
0 \to K^\wedge \to P^\wedge \to M^\wedge \to 0
$$
is a split exact sequence.
\end{lemma}

\begin{proof}
As $M$ is flat, each of the sequences
$$
0 \to K/I^nK \to P/I^nP \to M/I^nM \to 0
$$
is short exact, see
Lemma \ref{lemma-flat-tor-zero}
and the sequence $0 \to K^\wedge \to P^\wedge \to M^\wedge \to 0$
is a short exact sequence, see
Lemma \ref{lemma-completion-generalities}.
It suffices to show that we can find splittings
$s_n : M/I^nM \to P/I^nP$ such that $s_{n + 1} \bmod I^n = s_n$.
We will construct these $s_n$ by induction on $n$.
Pick any splitting $s_1$, which exists as $M/IM$ is a projective $R/I$-module.
Assume given $s_n$ for some $n > 0$. Set
$P_{n + 1} = \{x \in P \mid x \bmod I^nP \in \Im(s_n)\}$.
The map $\pi : P_{n + 1}/I^{n + 1}P_{n + 1} \to M/I^{n + 1}M$ is surjective
(details omitted). As $M/I^{n + 1}M$ is projective as a $R/I^{n + 1}$-module
by
Lemma \ref{lemma-lift-projective}
we may choose a section
$t : M/I^{n + 1}M \to P_{n + 1}/I^{n + 1}P_{n + 1}$
of $\pi$. Setting $s_{n + 1}$ equal to the composition of
$t$ with the canonical map $P_{n + 1}/I^{n + 1}P_{n + 1} \to P/I^{n + 1}P$
works.
\end{proof}

\begin{lemma}
\label{lemma-complete-modulo-nilpotent}
Let $A$ be a Noetherian ring. Let $I, J \subset A$ be ideals.
If $A$ is $I$-adically complete and $A/I$ is $J$-adically complete,
then $A$ is $J$-adically complete.
\end{lemma}

\begin{proof}
Let $B$ be the $(I + J)$-adic completion of $A$. By
Lemma \ref{lemma-completion-flat} $B/IB$ is the $J$-adic completion of $A/I$
hence isomorphic to $A/I$ by assumption. Moreover $B$ is $I$-adically
complete by Lemma \ref{lemma-complete-by-sub}. Hence $B$ is a finite
$A$-module by Lemma \ref{lemma-finite-over-complete-ring}.
By Nakayama's lemma (Lemma \ref{lemma-NAK} using $I$ is in the
Jacobson radical of $A$
by Lemma \ref{lemma-radical-completion}) we find that $A \to B$ is surjective.
The map $A \to B$ is flat by Lemma \ref{lemma-completion-flat}.
The image of $\Spec(B) \to \Spec(A)$ contains $V(I)$ and as $I$
is contained in the Jacobson radical of $A$ we find $A \to B$ is faithfully flat
(Lemma \ref{lemma-ff-rings}). Thus $A \to B$ is injective. Thus $A$ is
complete with respect to $I + J$, hence a fortiori complete
with respect to $J$.
\end{proof}









